% ============ 第4章 §4.9 Bochner-Riesz平均与Nikodym极大函数(原书 pp.371-378) ============
\section{Bochner--Riesz 平均与 Nikodym 极大函数}\label{sec:4.9}
\index{Bochner--Riesz 平均}
\index{Fefferman 定理}
\index{乘子!乘子算子(Rn 上的函数的)}


本节利用 Nikodym 极大函数的估计,讨论与 Fourier 反演公式相关的 Bochner--Riesz 平均问题.

正如在 4.3 节所确认的,对 $f \in \mathcal{S}(\mathbb{R}^{n})$,Fourier 反演公式
\[
\int_{\mathbb{R}^{n}}e^{2\pi i x\cdot\xi}\hat{f}(\xi)\dd\xi = f(x)
\]
成立.若 $f$ 是 $L^{p}(\mathbb{R}^{n})$ 中更一般的函数,则超函数意义下的 Fourier 反演公式成立,但问题在于了解更定量的收敛.于是,我们考察部分和
\[
S_{R}f(x) := \int_{|\xi| \leq R}e^{2\pi i x\cdot\xi}\hat{f}(\xi)\dd\xi \quad (R > 0)
\]
当 $R \to \infty$ 时在 $L^{p}(\mathbb{R}^{n})$ 中的收敛问题.

由一致有界性原理与扩张变换的不变性,该问题归结为取 $R = 1$ 时如下不等式是否成立:
\begin{equation}\tag{4.9.1}
\|S_{1}f\|_{L^{p}(\mathbb{R}^{n})} \leq C\|f\|_{L^{p}(\mathbb{R}^{n})}.
\label{eq:4.9.1}\end{equation}

\eqref{eq:4.9.1}在一维情形,由 Hilbert 变换的 $L^{p}(\mathbb{R})$ 有界性可知在 $1 < p < \infty$ 时成立.但是,在二维以上,只知道由 Plancherel 定理得到的 $p = 2$ 时成立.$S_{1}$ 的核的绝对值在远方以 $|x|^{-(n+1)/2}$ 衰减,因此考虑到对偶性,曾期望在
\[
p \in \left(\frac{2n}{n+1}, \frac{2n}{n-1}\right)
\]
时有界.然而,1971 年,C. Fefferman 证明了如下著名定理,表明该猜想不成立.

\begin{theorem}\label{thm:4.9.1}
在 $n > 1$,$p \neq 2$ 时\eqref{eq:4.9.1}不成立.
\end{theorem}

C. Fefferman 的定理是利用二维 Kakeya 集的 Minkowski 维数为 2 的反例证明的.这里略去证明,继续前进.

在 Fefferman 的反例出现以后,兴趣的对象转移到证明比 $S_{1}$ 奇异性更小的 Bochner--Riesz 平均
\[
S_{1}^{\lambda}f(x) := \int_{\mathbb{R}^{n}}e^{2\pi i x\cdot\xi}(1-|\xi|^{2})_{+}^{\lambda}\hat{f}(\xi)\dd\xi, \quad t_{+}^{\lambda} := \max(t,0)^{\lambda} \quad (\lambda > 0)
\]
的 $L^{p}(\mathbb{R}^{n})$ 有界性问题.由驻相法可以证明 $S_{1}^{\lambda}$ 的核的绝对值在远方以 $|x|^{-(n+1)/2-\lambda}$ 衰减,并且猜想在
\[
p \in \left[\frac{2n}{n+1}, \frac{2n}{n-1}\right], \quad \lambda > 0
\]
时 $S_{1}^{\lambda}$ 在 $L^{p}(\mathbb{R}^{n})$ 上有界.

更精确地,已经得到 Bochner--Riesz 平均在 $L^{p}(\mathbb{R}^{n})$ 上有界的必要条件.即,令
\[
\lambda(p) := \max\left\{n\left|\frac{1}{2} - \frac{1}{p}\right| - \frac{1}{2}, 0\right\},
\]
则在 $p \in [1, \infty]$ 时 $S_{1}^{\lambda}$ 在 $L^{p}(\mathbb{R}^{n})$ 上有界的必要条件为
\[
\lambda > \lambda(p).
\]
关于 Bochner--Riesz 平均的猜想是:这个必要条件也是充分条件.下面确认该猜想在二维时正确.在三维以上,只知道部分结果,仍悬而未决.

\textbf{C\'ordoba 对二维情形的解法.} 1977 年 C\'ordoba 用二维 Nikodym 极大函数的 $L^{2}(\mathbb{R}^{2})$ 估计,给出了 Bochner--Riesz 平均猜想在二维成立的另一证明.该问题最初的证明出自 Carleson 与 Sj\"olin,用的是驻相法.

这里我们概观一下 C\'ordoba 的证明.

以下,把由函数 $\phi$ 给出的 Fourier 乘子算子记为 $D(\phi)$.即
\[
D(\phi)f(x) := \int_{\mathbb{R}^{n}}e^{2\pi i x\cdot\xi}\phi(\xi)\hat{f}(\xi)\dd\xi.
\]
在 $p \in [4/3, 4]$,$\lambda > 0$ 时,目标是确认
\[
\|S_{1}^{\lambda}f\|_{L^{p}(\mathbb{R}^{2})} \leq C\|f\|_{L^{p}(\mathbb{R}^{2})}.
\]
考虑到对偶性以及与由 Plancherel 定理推出的 $L^{2}(\mathbb{R}^{2})$ 有界性之间的插值,只需证明
\begin{equation}\tag{4.9.2}
\|S_{1}^{\lambda}f\|_{L^{4}(\mathbb{R}^{2})} \leq C\|f\|_{L^{4}(\mathbb{R}^{2})}
\label{eq:4.9.2}\end{equation}
即可.

设 $\epsilon \in (0, 1/2]$.取满足
\[
\operatorname{supp}\phi_{\epsilon} \subset (1-\epsilon, 1-\epsilon/2), \quad |(d^{k}/dt^{k})\phi_{\epsilon}| \leq C_{k}\epsilon^{-k} \quad (k = 1, 2, \dots)
\]
的光滑实值函数 $\phi_{\epsilon}$,定义
\[
M_{\epsilon} := D\left(\phi_{\epsilon}(|\xi|)\right).
\]
此时,用单位分解把 $(1-|\xi|^{2})_{+}^{\lambda}$ 的支集作二进制分解,可得展开
\[
S_{1}^{\lambda}f(x) = M_{0}f(x) + \sum_{j=1}^{\infty}2^{-\lambda j}M_{2^{-j}}f(x).
\]
这里,$M_{0}$ 是以 $B(0, 1/2)$ 为支集的光滑函数给出的 Fourier 乘子算子,在 $L^{4}(\mathbb{R}^{2})$ 上有界.

现在,若对充分小的 $\epsilon > 0$ 一致地得到估计
\begin{equation}\tag{4.9.3}
\|M_{\epsilon}f\|_{L^{4}(\mathbb{R}^{2})} \leq C\epsilon^{-\lambda/2}\|f\|_{L^{4}(\mathbb{R}^{2})},
\label{eq:4.9.3}\end{equation}
则由上面的展开式显然可得\eqref{eq:4.9.2}.把 $\phi_{\epsilon}(|\xi|)$ 的支集
\[
\{\xi \in \mathbb{R}^{2}:\ 1-\epsilon < |\xi| < 1-\epsilon/2\}
\]
在角方向用长度 $\sim\sqrt{\epsilon}$ 的弯曲长方形 $\tilde{R}_{\epsilon,i}$($i = 1, 2, \ldots$)分解,把 $M_{\epsilon}$ 进一步展开为
\[
M_{\epsilon} = \sum_{i=1}^{\sim\epsilon^{-1/2}}D(\psi_{\epsilon,i}).
\]
这里,$\psi_{\epsilon,i}$ 是以 $\tilde{R}_{\epsilon,i}$ 为支集的光滑函数.

把通过各 $\tilde{R}_{\epsilon,i}$ 中心的单位向量记为 $\omega_{i} \in S^{1}$,不失一般性,以下假定 $\omega_{i}$ 的第一坐标至少为 $1/\sqrt{2}$ 来进行证明.另外,把这样的 $i$ 全体的集合记为 $I$,为简单起见虽稍显粗糙,但设
\[
M_{\epsilon} = \sum_{i \in I}D(\psi_{\epsilon,i}).
\]

由 Fourier 变换的基本公式,
\begin{equation}\tag{4.9.4}
\begin{aligned}
\|M_{\epsilon}f\|_{4}^{4} & = \int_{\mathbb{R}^{2}}\left|\sum_{i \in I}D(\psi_{\epsilon,i})f\right|^{4}\dd x \\
& = \int_{\mathbb{R}^{2}}\left|\sum_{i,j \in I}D(\psi_{\epsilon,i})f\,D(\psi_{\epsilon,j})f\right|^{2}\dd x \\
& = \int_{\mathbb{R}^{2}}\left|\sum_{i,j \in I}\left(D(\psi_{\epsilon,i})f\,D(\psi_{\epsilon,j})f\right)^{\wedge}\right|^{2}\dd\xi \\
& = \int_{\mathbb{R}^{2}}\left|\sum_{i,j \in I}\left(D(\psi_{\epsilon,i})f\right)^{\wedge}*\left(D(\psi_{\epsilon,j})f\right)^{\wedge}\right|^{2}\dd\xi.
\end{aligned}
\label{eq:4.9.4}\end{equation}
为估计该右端,准备如下几何引理.

\begin{lemma}\label{lem:4.9.2}
以 $c$ 为不依赖于 $\epsilon$ 的常数,
\[
\left\|\sum_{i,j \in I}\chi_{\tilde{R}_{\epsilon,i}} * \chi_{\tilde{R}_{\epsilon,j}}\right\|_{\infty} \leq c
\]
成立.
\end{lemma}

\begin{proof}
该引理本质上由如下事实推出:若 $\omega_{1}, \omega_{2}, \omega'_{1}, \omega'_{2} \in S^{1}$ 满足
\begin{align*}
0 &\neq \omega_{1} + \omega_{2} = \omega'_{1} + \omega'_{2}, \\
\intertext{则要么 $\omega_{1} = \omega'_{1}, \omega_{2} = \omega'_{2}$,要么 $\omega_{1} = \omega'_{2}, \omega_{2} = \omega'_{1}$. 设}
g(\xi) &:= \sum_{i,j \in I}\left(\chi_{\tilde{R}_{\epsilon,i}} * \chi_{\tilde{R}_{\epsilon,j}}\right)(\xi).
\end{align*}
若取 $g$ 的支集上的点为 $\zeta_{0}$,则存在 $\xi_{0} \in \tilde{R}_{\epsilon,i_{0}}$,$\eta_{0} \in \tilde{R}_{\epsilon,j_{0}}$ 使得
\[
\zeta_{0} = \xi_{0} + \eta_{0}.
\]
现在设存在 $\xi_{1} \in \tilde{R}_{\epsilon,i_{1}}$,$\eta_{1} \in \tilde{R}_{\epsilon,j_{1}}$ 使得
\begin{align*}
\zeta_{0} &= \xi_{1} + \eta_{1}. \\
\intertext{令}
\zeta'_{0} &:= \frac{\xi_{0}}{|\xi_{0}|} + \frac{\eta_{0}}{|\eta_{0}|}, \quad \zeta'_{1} := \frac{\xi_{1}}{|\xi_{1}|} + \frac{\eta_{1}}{|\eta_{1}|}, \\
\intertext{则}
&\begin{array}{r l}
|\zeta'_{0} - \zeta'_{1}| & \leq |\zeta'_{0} - \zeta_{0}| + |\zeta_{0} - \zeta'_{1}| \\
& = \left|\frac{\xi_{0}}{|\xi_{0}|} - \xi_{0}\right| + \left|\frac{\eta_{0}}{|\eta_{0}|} - \eta_{0}\right| + \left|\frac{\xi_{1}}{|\xi_{1}|} - \xi_{1}\right| + \left|\frac{\eta_{1}}{|\eta_{1}|} - \eta_{1}\right| \\
& \leq 4\epsilon.
\end{array}
\end{align*}
由此,满足上述条件的 $\xi_{1}, \eta_{1}$ 的个数不多是显然的.特别地,对 $\zeta'_{0} \sim (2, 0)$ 的情形,用 $1 - \cos\theta \sim \theta^{2}/2$($\theta \sim 0$)的关系考察,可知其个数不依赖于 $\epsilon$.
\end{proof}

用该引理估计\eqref{eq:4.9.4},得
\begin{equation}\tag{4.9.5}
\|M_{\epsilon}f\|_{4}^{4} \leq C\int_{\mathbb{R}^{2}}\left|\sum_{i \in I}|D(\psi_{\epsilon,i})f|^{2}\right|^{2}\dd x
\label{eq:4.9.5}\end{equation}
成立.

来估计 $D(\psi_{\epsilon,i})f(x)$.改写
\[
D(\psi_{\epsilon,i})f(x) = \iint e^{2\pi i(x-y)\cdot\xi}\psi_{\epsilon,i}(\xi)\dd\xi f(y)\dd y =: \int K_{i}(x-y)f(y)\dd y,
\]
来估计
\begin{equation}\tag{4.9.6}
K_{i}(x) = \int e^{2\pi i x\cdot\xi}\psi_{\epsilon,i}(\xi)\dd\xi.
\label{eq:4.9.6}\end{equation}

对每个 $\omega_{i}$,各选取一个满足 $U_{i}\omega_{i} = (1, 0)$ 的 $\mathbb{R}^{2}$ 的直交变换 $U_{i}$,令
\begin{align*}
R_{\epsilon} &:= \left\{x \in \mathbb{R}^{2}:\ |x_{1}| \leq \frac{1}{\epsilon},\ |x_{2}| \leq \frac{1}{\sqrt{\epsilon}}\right\}, \\
\intertext{设}
R_{\epsilon,i} &:= U_{i}^{-1}R_{\epsilon}.
\end{align*}
此时,可以证明只需证明
\begin{equation}\tag{4.9.7}
|K_{i}(x)| \leq C\frac{\chi_{R_{\epsilon,i}}(x)}{|R_{\epsilon,i}|}
\label{eq:4.9.7}\end{equation}
即可.这一点留到最后确认,先继续前进.

对每个 $\tilde{R}_{\epsilon,i}$($i \in I$),选取满足 $\tilde{R}_{\epsilon,i} \subset Q_{\epsilon,i}$ 的边长为 $\sim\sqrt{\epsilon}$ 的正方形 $Q_{\epsilon,i}$,令
\begin{align*}
f_{i} &:= D\left(\chi_{Q_{\epsilon,i}}\right)f, \\
\intertext{则显然}
D(\psi_{\epsilon,i})f &= D(\psi_{\epsilon,i})f_{i}.
\end{align*}
由该式与\eqref{eq:4.9.5},\eqref{eq:4.9.7}得
\begin{equation}\tag{4.9.8}
\|M_{\epsilon}f\|_{4}^{2} \leq C\left\|\sum_{i \in I}\left(\frac{\chi_{R_{\epsilon,i}}}{|R_{\epsilon,i}|} * |f_{i}|\right)^{2}\right\|_{2}.
\label{eq:4.9.8}\end{equation}
用对偶性来解析该右端.

任取非负函数 $w \in L^{2}(\mathbb{R}^{2})$,$\|w\|_{2} = 1$.利用 $\left\|\frac{\chi_{R_{\epsilon,i}}}{|R_{\epsilon,i}|}\right\|_{1} = 1$,
\begin{align*}
&\begin{array}{r l}
\int\left(\frac{\chi_{R_{\epsilon,i}}}{|R_{\epsilon,i}|} * |f_{i}|\right)^{2}w\dd x & \leq \int\left(\frac{\chi_{R_{\epsilon,i}}}{|R_{\epsilon,i}|} * |f_{i}|^{2}\right)w\dd x \\
& = \int|f_{i}|^{2}\left(\frac{\chi_{R_{\epsilon,i}}}{|R_{\epsilon,i}|} * w\right)\dd x.
\end{array} \\
\intertext{用此式,}
&\begin{array}{r l}
\int\sum_{i \in I}\left(\frac{\chi_{R_{\epsilon,i}}}{|R_{\epsilon,i}|} * |f_{i}|\right)^{2}w\dd x & \leq \int\left(\sum_{i \in I}|f_{i}|^{2}\right)\left(\sup_{i \in I}\frac{\chi_{R_{\epsilon,i}}}{|R_{\epsilon,i}|} * w\right)\dd x \\
& \leq \int\left(\sum_{i \in I}|f_{i}|^{2}\right)K_{N,a}w\dd x \\
& \leq \left\|\sum_{i}|f_{i}|^{2}\right\|_{2}\cdot\|K_{N,a}w\|_{2}.
\end{array}
\end{align*}
这里,取 $N = a = \frac{1}{\sqrt{\epsilon}}$,$K_{N,a}w$ 表示 4.7 节定义的 Nikodym 极大函数.

这里承认(不加证明地使用)如下事实:
\begin{align*}
\left\|\sum_{i}|f_{i}|^{2}\right\|_{2} &\leq C\|f\|_{4}^{2}, \\
\intertext{因此问题归结为 Nikodym 极大函数 $K_{N,a}w$ 的 $L^{2}(\mathbb{R}^{2})$ 估计.正如在 4.7 节所确认的,它可估计为}
\|K_{N,a}w\|_{2} &\leq C(\log N)^{1/2}\|w\|_{2} = C(\log N)^{1/2},
\end{align*}
由\eqref{eq:4.9.8}得
\[
\|M_{\epsilon}f\|_{4} \leq C(\log(1/\epsilon))^{1/4}\|f\|_{4}.
\]
由此\eqref{eq:4.9.3}成立.

\textbf{验证\eqref{eq:4.9.7}.} 设 $\psi_{0} \in C_{0}^{\infty}(\mathbb{R}^{2})$ 满足
\begin{align*}
\operatorname{supp}\psi_{0} &\subset \left\{\xi = (\xi_{1}, \xi_{2}) \in \mathbb{R}^{2}:\ 1-\epsilon < |\xi| < 1-\epsilon/2,\ -\sqrt{\epsilon} < \xi_{2} < \sqrt{\epsilon}\right\}, \\
\intertext{且}
\left|(\partial/\partial\xi_{1})^{\alpha_{1}}(\partial/\partial\xi_{2})^{\alpha_{2}}\psi_{0}\right| &\leq C_{\alpha_{1},\alpha_{2}}\epsilon^{-\alpha_{1}}\epsilon^{-\alpha_{2}/2}. \\
\intertext{不失一般性,可以设}
\psi_{\epsilon,i}(\xi) &= \psi_{0}(U_{i}\xi). \\
\intertext{由\eqref{eq:4.9.6},}
&\begin{array}{r l}
K_{i}(x) & = \int e^{2\pi i x\cdot\xi}\psi_{\epsilon,i}(\xi)\dd\xi \\
& = \int e^{2\pi i U_{i}^{-1}x\cdot U_{i}\xi}\psi_{0}(U_{i}\xi)\dd\xi \\
& = \int e^{2\pi i U_{i}^{-1}x\cdot\xi}\psi_{0}(\xi)\dd\xi.
\end{array} \\
\intertext{因此,为证明主张,只需考察}
K_{0}(x) &:= \int e^{2\pi i x\cdot\xi}\psi_{0}(\xi)\dd\xi. \\
\intertext{利用对 $\psi_{0}$ 的假设作分部积分,可知}
(\epsilon|x_{1}|)^{m}(\sqrt{\epsilon}|x_{2}|)^{m}|K_{0}(x)| &\leq C_{m}\epsilon^{3/2} \quad (m = 0, 1, \dots)
\end{align*}
成立.由此计算出
\[
|K_{0}(x)| \leq C_{m}\sum_{k=0}^{\infty}2^{-(m-2)k}\frac{\chi_{2^{k}R_{\epsilon}}}{|2^{k}R_{\epsilon}|},
\]
由 Minkowski 不等式使主张得到验证.
